\documentclass{article}

\usepackage{amsmath,amssymb}
\usepackage{xurl}
\usepackage[hidelinks]{hyperref}

% Shared commands for notes in docs/paper-gaps/.

\DeclareUnicodeCharacter{2080}{\ifmmode{}_{0}\else\textsubscript{0}\fi}
\DeclareUnicodeCharacter{2081}{\ifmmode{}_{1}\else\textsubscript{1}\fi}
\DeclareUnicodeCharacter{2082}{\ifmmode{}_{2}\else\textsubscript{2}\fi}
\DeclareUnicodeCharacter{2099}{\ifmmode{}_{n}\else\textsubscript{n}\fi}

\newcommand{\Lean}{\textsc{Lean}}
\ExplSyntaxOn
\NewDocumentCommand{\leanid}{m}{
  \ifmmode
    \textsf{\detokenize{#1}}
  \else
    \group_begin:
    \urlstyle{sf}
    \tl_set:Nn \l_tmpa_tl {#1}
    \tl_replace_all:Nnn \l_tmpa_tl {\_} {_}
    \exp_args:NV \path \l_tmpa_tl
    \group_end:
  \fi
}
\ExplSyntaxOff
\providecommand{\tr}{\operatorname{tr}}
\newcommand{\tnleanIssueBase}{https://github.com/LionSR/TNLean/issues/}
\newcommand{\tnleanPrBase}{https://github.com/LionSR/TNLean/pull/}
\newcommand{\tnleanDocBase}{https://sirui-lu.com/TNLean/docs/}
\newcommand{\ghissue}[1]{\href{\tnleanIssueBase#1}{GitHub issue \##1}}
\newcommand{\ghpr}[1]{\href{\tnleanPrBase#1}{PR \##1}}
\newcommand{\doclink}[1]{\href{\tnleanDocBase#1.html}{\path{#1}}}

% Dirac notation and the identity operator, as in blueprint/src/macros/common.tex.
% \providecommand keeps note-local definitions authoritative.
\usepackage{dsfont}
\providecommand{\ket}[1]{|#1\rangle}
\providecommand{\bra}[1]{\langle #1|}
\providecommand{\braket}[2]{\langle #1 | #2 \rangle}
\providecommand{\ketbra}[2]{|#1\rangle\!\langle #2|}
\providecommand{\Id}{\mathds{1}}
\providecommand{\one}{\mathds{1}}
\providecommand{\C}{\mathbb{C}}
\providecommand{\MN}[1]{M_{#1}(\C)}
\providecommand{\MD}{M_D(\C)}

% External referencing for paper sources.  The build helper writes sanitized
% auxiliary files under build/paper-aux/ and excludes duplicated source labels.
\usepackage{xr}

% Import a generated paper auxiliary file with a caller-supplied label prefix.
% The candidate paths support builds from docs/paper-gaps/, the repository root,
% and build/paper-aux/ itself.
\newcommand{\importpaperaux}[2]{%
  \IfFileExists{../../build/paper-aux/#2.aux}{%
    \externaldocument[#1]{../../build/paper-aux/#2}%
  }{%
    \IfFileExists{build/paper-aux/#2.aux}{%
      \externaldocument[#1]{build/paper-aux/#2}%
    }{%
      \IfFileExists{#2.aux}{%
        \externaldocument[#1]{#2}%
      }{}%
    }%
  }%
}

% General external reference macros
\newcommand{\paperref}[2]{\ref{#1-#2}}
\newcommand{\papereqref}[2]{(\ref{#1-#2})}

% CPSV16 source (1606.00608 / MPDO-22-12-17-2.tex) external document and shortcuts
\importpaperaux{cpsv16-}{1606.00608/MPDO-22-12-17-2-tnlean}
\newcommand{\cpsvref}[1]{\paperref{cpsv16}{#1}}
\newcommand{\cpsveqref}[1]{\papereqref{cpsv16}{#1}}

% Verdict marker: \gapnote{<kind>}{<status>}. Kind is the policy
% classification (clarification, local-correction, scope-restriction,
% unfaithful, false-source, open-gap); status is open, wip (elimination
% underway), resolved, or historical. Machine-read by the site index;
% prints a verdict line.
\newcommand{\gapnote}[2]{\par\noindent\textbf{Verdict:} #1 (#2).\par}


\title{Fixed-Length Chain Equality Does Not Imply Combined-Tensor MPV Equality}
\author{}
\date{2026-07-20}

\begin{document}
\maketitle

\section{The incompatible statements}

For a site-dependent closed chain
$\mathbf A=(A_0,\ldots,A_{n-1})$, equality of the fixed-length states tests
only cyclic products containing one matrix from each site, in their cyclic
order.  In contrast, MPV equality for the tensor obtained by combining the
site and physical indices tests every word in the enlarged alphabet.  Such a
word may repeat one site and omit every other site.  Therefore fixed-length
state equality cannot imply combined-tensor MPV equality.

This distinction is visible already for physical and bond dimensions one and
chain length three.  Take
\[
  A_0=A_1=A_2=1,
  \qquad
  B_0=2,\quad B_1=\tfrac12,\quad B_2=1.
\]
Every local tensor of both chains is injective, and the two closed-chain
coefficients agree because
$A_0A_1A_2=B_0B_1B_2=1$.  The combined tensors do not generate the same MPV:
the length-one word selecting site zero has coefficient $1$ for the combined
tensor of $\mathbf A$ and coefficient $2$ for that of $\mathbf B$.

Consequently the former assertion
\path{MPSChainTensor.SameStateBridgeHyp} has no construction even after adding
injectivity of the second chain.  It is not the blocking argument in
\cite[Theorem~\textup{thm:inj\_MPS}, lines~688--725]{Molnar2018NormalPEPS}.

\section{The source-faithful route}

The source assumes that every local tensor in both chains is injective and
derives one invertible gauge per bond directly from equality of the single
closed-chain state.  This statement is formalized as
\path{TNLean.PEPS.fundamentalTheorem_injectiveMPSChain_of_sameState} and is
recorded in Chapter~24 of the blueprint.  It does not pass through an
all-length MPV assertion for the combined tensor.

For constant chains, equality at one length likewise determines a tensor only
up to a scalar whose $n$-th power is one and a similarity transformation.  A
phase-free conclusion from one fixed length is false: in bond dimension one,
replace $A=1$ by $B=\omega$ for a nontrivial $n$-th root of unity $\omega$.
The length-$n$ states agree, but scalar conjugation cannot change $1$ into
$\omega$.  The all-length hypothesis used by
\path{MPSChainTensor.ti_reduction_corollary} rules out this phase and remains a
separate, stronger theorem.

\section{Resolution}

The abstract same-state reduction hypothesis and its consumers are removed.
The fixed-length injective-chain Fundamental Theorem is represented only by
the direct source-faithful theorem in Chapter~24.  The Chapter~23 results based
on combined-tensor MPV equality remain explicitly all-length results and are
not presented as consequences of fixed-length state equality.

\end{document}
